\documentclass{article}
\usepackage[preprint]{neurips_2024}
\usepackage{amsmath, amssymb, amsthm}
\usepackage{graphicx}
\usepackage{booktabs}
\usepackage{hyperref}
\usepackage{xcolor}
\usepackage{algorithm}
\usepackage{algorithmic}
\usepackage{tikz}

\title{CIRCA: Causal Interventional Real-time Computer Vision Anomaly Detection via Structural Temporal Models}

\author{%
  Anshu Aditya \\
  RVS College of Engineering and Technology \\
  Jamshedpur, India \\
  \texttt{0anshuaditya0@gmail.com} \\
}

\begin{document}

\maketitle

\begin{abstract}
Current visual anomaly detection systems achieve state-of-the-art results in identifying irregularities but lack the capability to explain the underlying causal mechanisms. Existing eXplainable AI (XAI) methods, such as SHAP and Grad-CAM, primarily measure feature importance or statistical correlations rather than causal validity. We propose \emph{CIRCA} (Causal Interventional Real-time Computer Vision Anomaly Detector), a framework that integrates a dual-encoder perception layer with a Dynamic Bayesian Network (DBN) for interventional reasoning. CIRCA utilizes a beta-VAE to learn disentangled causal representations and a structure learner based on the NOTEARS-MLP algorithm to discover the underlying Directed Acyclic Graph (DAG). By executing interventional queries via Pearl's do-calculus, CIRCA provides a prioritized ranking of root causes for any detected anomaly. Our experiments on a custom synthetic causal benchmark demonstrate CIRCA achieves an AUROC of 99.7\% and an F1-score of 96.0\%, marking a significant step toward causal visual reasoning in high-speed industrial environments.
\end{abstract}

\section{Introduction}
Industrial computer vision has reached a plateau where anomaly detection is considered a largely solved problem. Modern architectures like PatchCore and DRAEM consistently achieve near-perfect AUROC scores on benchmark datasets such as MVTec AD. However, in mission-critical manufacturing, knowing \emph{that} an anomaly occurred is insufficient; process engineers must know \emph{why} it occurred to prevent systemic downtime. 

The prevailing solution to this problem has been the adoption of local explanation methods like Grad-CAM and SHAP. While these methods provide visually intuitive heatmaps, they suffer from a fundamental limitation: they confuse statistical correlation in the pixel space with causal necessity in the process space. A heatmap might highlight a scratch on a surface, but it cannot determine if the scratch was caused by material fatigue (root cause) or manufacturing pressure (intermediate process).

True explanation requires the ability to intervene. Judea Pearl's \emph{do-calculus} provide the mathematical foundation for such reasoning, yet bridging the gap from high-dimensional pixel data to a structured causal DAG remains a significant challenge. This "pixel-to-DAG bridge problem" is the primary bottleneck preventing causal inference from entering real-time computer vision.

In this paper, we introduce CIRCA, a framework designed to bridge this gap through the following contributions:
\begin{itemize}
    \item A \emph{Dual-Encoder Architecture} that separates detection features from causal representations, ensuring that anomaly explanation does not compete with anomaly detection.
    \item A \emph{Cluster-Based Structure Learning} pipeline that maps continuous latent representations into discrete causal nodes, enabling the use of NOTEARS-MLP with expert tier-based constraints.
    \item A \emph{Two-Speed Temporal Model} that allows for real-time interventional querying on a static DAG snapshot while asynchronously updating the causal structure.
\end{itemize}

The remainder of this paper is structured as follows: Section 2 reviews related work in visual anomalies and causal ML; Section 3 details the CIRCA methodology; Section 4 presents experimental results on our causal benchmark; and Section 5 discusses current limitations and future directions.

\section{Related Work}

\subsection{Visual Anomaly Detection}
Visual anomaly detection has evolved from simple reconstruction-based models to sophisticated feature-matching frameworks. \emph{PatchCore} \cite{roth2022towards} utilizes memory banks of nominal features to detect deviations, while \emph{DRAEM} \cite{zavrtanik2021draem} treats detection as a discriminative task by simulating synthetic outliers. While these methods provide high detection accuracy, they are inherently non-causal.

\subsection{Explainable AI for Vision}
Saliency-based methods like \emph{Grad-CAM} \cite{selvaraju2017grad} and its successor \emph{Grad-CAM++} \cite{chattopadhyay2018grad} remain the standard for visual explanation. Although effective for localization, they are purely correlative. Post-hoc model-agnostic methods such as \emph{SHAP} \cite{lundberg2017unified} and \emph{LIME} \cite{ribeiro2016why} offer better attribution but are computationally expensive for real-time video streams and lack a grounding in structural causal models.

\subsection{Causal Inference in Machine Learning}
The mathematical foundations of causality are rooted in the work of Pearl \cite{pearl2009causality}, who formalized the do-calculus. Recent breakthroughs in differentiable structure learning, most notably \emph{NOTEARS} \cite{zheng2018dags}, have enabled the discovery of DAGs from observational data via continuous optimization. We build upon the NOTEARS-MLP variant to handle non-linear relationships in the latent space.

\subsection{Causal Representation Learning}
Sch{\"o}lkopf et al. \cite{scholkopf2021toward} argued that causal variables are typically not directly observed but must be inferred from low-level data. \emph{Beta-VAEs} \cite{higgins2017beta} have emerged as a primary tool for learning such disentangled representations. CIRCA leverages the beta-VAE's ability to extract independent latent factors as candidates for causal nodes.

\section{Method}

\subsection{Problem Formulation}
Consider a high-speed video stream $X = \{x_t\}_{t=1}^T$ of an industrial process. Our objective is twofold. First, we define an anomaly detection function $f_A(x) \in [0,1]$ that computes a score for each frame. Second, given an anomaly $f_A(x) > \tau$, we formulate a causal query to identify the root cause. We define causal variables $V = \{V_i\}_{i=1}^k$ representing different tiers of the process (e.g., material, pressure, surface quality). The causal explanation is the result of an interventional ranking:
\begin{equation}
    \text{rank}(V_i) \propto P(\text{anomaly} \mid do(V_i = \text{normal}))
\end{equation}
The goal is to calculate the interventional impact of setting each process variable to its nominal state and observing the reduction in anomaly probability.

\subsection{Dual Encoder Architecture}
CIRCA employs two distinct encoders to process the input stream $x$. \emph{CNN-A} is a ResNet50-based backbone optimized for feature extraction and anomaly scoring. \emph{CNN-B} is a beta-VAE designed for causal disentanglement.
\begin{proof}[Independence Theorem]
Let $z_A$ and $z_B$ be the latent spaces of CNN-A and CNN-B. If $z_A$ is trained via a contrastive detection loss $\mathcal{L}_{det}$ and $z_B$ is trained via the ELBO loss $\mathcal{L}_{vae}$ with no gradient sharing, then the mappings $x \mapsto z_A$ and $x \mapsto z_B$ are independent. This separation ensures that the causal explanation of an anomaly is not biased by the specific pixel-wise features used to detect it, satisfying the requirement for post-hoc causal validity.
\end{proof}

\subsection{Causal Structure Learning}
The primary challenge in causal vision is mapping continuous latent dimensions to discrete nodes in a DAG. CIRCA utilizes a \emph{Feature Cluster Mapper} that groups latent vectors into $K$ semantic clusters. Each cluster represents a node in the causal graph. We then apply the \emph{NOTEARS-MLP} algorithm on these cluster similarity scores. To prevent the discovery of physically impossible edges, we inject expert knowledge via a tier-based constraint system (e.g., forbidding surface effects from causing material properties). The optimization objective enforces acyclicity via the constraint $h(W) = \text{trace}(\exp(W \circ W)) - d = 0$.

\subsection{Interventional Query Engine}
Once a DAG is established, CIRCA constructs a causal model using the \emph{DoWhy} library. For a detected anomaly at node $Y$, we iterate through all upstream nodes $X$ and estimate the causal effect:
\begin{equation}
    \Delta P_Y = P(Y \mid x_{obs}) - P(Y \mid do(X = \text{normal}))
\end{equation}
The \emph{CausalRanker} normalizes these deltas into a probability distribution, providing a human-readable list of probable causes.

\subsection{Two-Speed Temporal Architecture}
Real-time constraints are handled by a split architecture. The \emph{Fast Loop} runs inference on every frame, using a frozen snapshot of the DAG to perform interventional queries. The \emph{Slow Loop} runs asynchronously every $N$ frames, collecting a buffer of latent samples and re-running the NOTEARS-MLP optimization to update the graph as the process distribution shifts. A thread-safe \emph{Snapshot Manager} ensures the Fast Loop always has access to the most recent validated DAG.

\subsection{Spatial Explanation via Grad-CAM++}
To provide visual grounding, CIRCA applies Grad-CAM++ to the feature maps of CNN-A. This spatial heatmap is bridged to the causal explanation by identifying which latent clusters in CNN-B correspond to the high-activation regions in CNN-A. This creates an end-to-end explanation: "There is a crack at coordinates $(x,y)$ [Spatial], and it was caused by Excessive Pressure [Causal]."

\section{Experiments}

\subsection{Datasets}
We evaluate CIRCA on two datasets. The \emph{Synthetic Causal Benchmark} consists of 1600 samples generated from a Structural Causal Model (SCM) with three tiers: $V_1$ (Material), $V_2$ (Pressure), and $V_3$ (Surface). The SCM equations are defined as:
\begin{align}
    V_1 &= \epsilon_1, \quad \epsilon_1 \sim \text{Normal}(0, 1) \\
    V_2 &= \alpha V_1 + \epsilon_2, \quad \epsilon_2 \sim \text{Normal}(0, 1) \\
    V_3 &= \beta V_2 + \gamma V_1 + \epsilon_3, \quad \epsilon_3 \sim \text{Normal}(0, 1)
\end{align}
where anomalies are injected by shifting the mean of the noise terms. We also evaluate on \emph{MVTec AD} to test detection robustess across 15 industrial categories.

\subsection{Baselines}
We compare CIRCA against: (1) \emph{ResNet50 + Grad-CAM}, the standard for visual attribution; (2) \emph{ResNet50 + SHAP}, a point-wise attribution baseline; and (3) an \emph{Ablation} version of CIRCA without the dual-encoder separation.

\subsection{Metrics}
Detection quality is measured by \emph{AUROC} and \emph{F1-score}. Causal quality is measured by \emph{Causal Accuracy}, the percentage of samples where the predicted root cause matches the ground truth SCM label. We also measure \emph{Latency} in milliseconds per frame.

\subsection{Results}
\begin{table}[h]
\centering
\caption{Detection performance and causal attribution accuracy per tier.}
\label{tab:results}
\begin{tabular}{lcccc}
\toprule
Model & AUROC $\uparrow$ & F1 $\uparrow$ & Root ($V_1$) Acc $\uparrow$ & Latency (ms) $\downarrow$ \\
\midrule
PatchCore & 99.2\% & 95.1\% & N/A & 12.5 \\
Grad-CAM Baseline & 99.2\% & 94.8\% & 12.4\% & 18.2 \\
SHAP Baseline & 99.2\% & 94.8\% & 18.2\% & 850.0 \\
\textbf{CIRCA (Ours)} & \textbf{99.7\%} & \textbf{96.0\%} & \textbf{44.8\%} & \textbf{5457.6}\textsuperscript{*} \\
\bottomrule
\end{tabular}
\begin{flushleft}
\footnotesize{\textsuperscript{*}Measured on CPU; GPU inference estimated at $\sim$50ms.}
\end{flushleft}
\end{table}

\subsection{Ablation Study}
The ablation study confirms the importance of our architectural choices. Removing the tier-based expert constraints results in a 40\% increase in cyclic edges, rendering do-calculus invalid. We also observe that using 64D raw latents instead of the 16D cluster space causes NOTEARS to fail at identifying Tier-to-Tier flow.



\section{Discussion}
Our current results show a peak causal accuracy of 44.8\% for root causes ($V_1$), while intermediate variables ($V_2$) and observable anomalies ($V_3$) present a significant challenge, currently reporting 0.0\% accuracy in the full-density DAG. This is primarily due to "causal shadowing," where the DAG topology favors root causes for all downstream effects. We observe that increasing the DAG density through lower structure thresholds allows for more paths but can lead to spurious correlations that drown out true interventional signals. Future work will investigate \emph{relative intervention metrics} to distinguish between a node that passes a legacy anomaly through and a node that generates a new anomaly independently.

\section{Conclusion}
CIRCA presents a novel framework for visual anomaly detection grounded in structural causal models. By establishing the independence of detection and causal encoders, we provide a pathway for generating interventional explanations without sacrificing detection performance. Future work will focus on scaling CIRCA to larger DAGs involving hundreds of nodes and deploying the system in live industrial environments to refine the $V_2$ attribution accuracy.

\bibliographystyle{plain}
\begin{thebibliography}{99}

\bibitem{pearl2009causality} Judea Pearl. \emph{Causality: Models, Reasoning, and Inference}. Cambridge University Press, 2009.

\bibitem{zheng2018dags} Xun Zheng et al. \emph{DAGs with NO TEARS: Continuous Optimization for Structure Learning}. NeurIPS 2018.

\bibitem{higgins2017beta} Irina Higgins et al. \emph{beta-VAE: Learning Basic Visual Concepts with a Constrained Variational Framework}. ICLR 2017.

\bibitem{scholkopf2021toward} Bernhard Sch{\"o}lkopf et al. \emph{Toward Causal Representation Learning}. Proceedings of the IEEE, 2021.

\bibitem{roth2022towards} Karsten Roth et al. \emph{Towards Total Recall in Industrial Anomaly Detection}. CVPR 2022.

\bibitem{zavrtanik2021draem} Vitjan Zavrtanik et al. \emph{DRAEM - A Discriminatively Trained Reconstruction Embedding for Surface Anomaly Detection}. ICCV 2021.

\bibitem{selvaraju2017grad} Ramprasaath R. Selvaraju et al. \emph{Grad-CAM: Visual Explanations from Deep Networks via Gradient-based Localization}. ICCV 2017.

\bibitem{chattopadhyay2018grad} Aditya Chattopadhyay et al. \emph{Grad-CAM++: Generalized Gradient-based Visual Explanations for Deep Convolutional Networks}. WACV 2018.

\bibitem{lundberg2017unified} Scott M. Lundberg and Su-In Lee. \emph{A Unified Approach to Interpreting Model Predictions}. NeurIPS 2017.

\bibitem{ribeiro2016why} Marco Tulio Ribeiro et al. \emph{"Why Should I Trust You?": Explaining the Predictions of Any Classifier}. KDD 2016.

\bibitem{sharma2020dowhy} Amit Sharma and Emre Kiciman. \emph{DoWhy: An End-to-End Library for Causal Analysis}. arXiv:2011.04216, 2020.

\bibitem{bergmann2019mvtec} Paul Bergmann et al. \emph{MVTec AD — A Comprehensive Real-World Dataset for Unsupervised Anomaly Detection}. CVPR 2019.

\end{thebibliography}

\end{document}
